\documentclass[12pt, a4paper]{article}

% ── Packages ──────────────────────────────────────────────────────────────────
\usepackage{amsmath, amssymb, amsthm}
\usepackage{geometry}
\usepackage{booktabs}
\usepackage{array}
\usepackage{hyperref}
\usepackage{xcolor}
\usepackage{tcolorbox}
\usepackage{enumitem}
\usepackage{parskip}

\geometry{margin=2.5cm}

% ══════════════════════════════════════════════════════════════════════════════
% ── Student-mode switch ───────────────────────────────────────────────────────
%
%   Compile with: \def\studentflag{1}\input{...}  → student version
%   Compile with: \input{...}                      → instructor version
%
\newif\ifstudentmode
\studentmodefalse       % default: instructor mode

\ifx\studentflag\undefined
  % no flag passed — keep studentmodefalse
\else
  \studentmodetrue
\fi
% ══════════════════════════════════════════════════════════════════════════════

% ── Student-mode macros ───────────────────────────────────────────────────────
%
%  \sol[<width>]{<math content>}   — blank inside a math environment
%  \soleq[<height>]{<math content>} — blank replacing an entire display equation
%
\newcommand{\sol}[2][6cm]{%
  \ifstudentmode
    \underline{\hspace{#1}}%
  \else
    {#2}%
  \fi
}

\newcommand{\soleq}[2][3cm]{%
  \ifstudentmode
    \par\vspace{4pt}%
    \noindent\rule{\linewidth}{0.5pt}\\[#1]%
    \noindent\rule{\linewidth}{0.5pt}%
    \par\vspace{4pt}%
  \else
    \[#2\]%
  \fi
}
% ──────────────────────────────────────────────────────────────────────────────

% ── Colors & boxes ────────────────────────────────────────────────────────────
\definecolor{mainblue}{RGB}{30, 100, 180}
\definecolor{boxbg}{RGB}{235, 243, 255}
\definecolor{notebg}{RGB}{255, 248, 230}

\tcbuselibrary{skins, breakable}

\newtcolorbox{resultbox}{
  colback=boxbg, colframe=mainblue,
  boxrule=1.2pt, arc=4pt,
  title={Result},
  fonttitle=\bfseries\color{white},
  coltitle=white, attach boxed title to top left={yshift=-2mm, xshift=4mm},
  boxed title style={colback=mainblue, arc=3pt}
}

\newtcolorbox{notebox}{
  colback=notebg, colframe=orange!70!black,
  boxrule=1pt, arc=4pt,
  left=4pt, right=4pt
}

% ── Theorem environments ───────────────────────────────────────────────────────
\theoremstyle{definition}
\newtheorem{example}{Example}[section]

% ── Math macros ───────────────────────────────────────────────────────────────
\newcommand{\bth}{\boldsymbol{\theta}}
\newcommand{\bx}{\mathbf{x}}
\newcommand{\bX}{\mathbf{X}}
\newcommand{\by}{\mathbf{y}}
\newcommand{\bw}{\mathbf{w}}
\newcommand{\yhat}{\hat{y}}
\newcommand{\xbar}{\bar{x}}
\newcommand{\ybar}{\bar{y}}
\newcommand{\sigmoid}{\sigma}

% ── Document ──────────────────────────────────────────────────────────────────
\begin{document}

\begin{center}
  {\LARGE \textbf{Section 1.2: Logistic Regression}} \\[0.5em]
  {\large AI-IM Class — Lecture Notes}%
  \ifstudentmode
    \quad {\normalsize\color{gray}[Student Copy]}%
  \fi
\end{center}

\vspace{1em}
\hrule
\vspace{1.5em}

% ══════════════════════════════════════════════════════════════════════════════
\section{From Regression to Classification}

In Section~1.1 we predicted a \emph{continuous} output $y \in \mathbb{R}$.
Now suppose the output is a \textbf{binary label}: $y \in \{0, 1\}$.
For example, we might ask:

\begin{center}
  \renewcommand{\arraystretch}{1.2}
  \begin{tabular}{ll}
    \toprule
    \textbf{Input $x$} & \textbf{Label $y$}        \\
    \midrule
    Tumor size (cm)    & 0 = benign, 1 = malignant \\
    Hours studied      & 0 = fail, 1 = pass        \\
    Email features     & 0 = not spam, 1 = spam    \\
    \bottomrule
  \end{tabular}
\end{center}

\subsection{A Toy Example}

Consider six observations, each labelled either class F or class T.
We encode the labels numerically as $y = 0$ for F and $y = 1$ for T:

\begin{center}
  \renewcommand{\arraystretch}{1.2}
  \begin{tabular}{c|cccccc}
    \toprule
    $x$   & 1 & 2 & 3 & 4 & 5 & 6 \\
    \midrule
    Class & F & F & T & F & T & T \\
    $y$   & 0 & 0 & 1 & 0 & 1 & 1 \\
    \bottomrule
  \end{tabular}
\end{center}

Small $x$ is mostly F and large $x$ is mostly T, but the classes
\emph{overlap} around $x = 3$ and $x = 4$.
This is the region where we should be \emph{uncertain} about the label.

\textbf{A naive attempt: linear regression.}
Since $y$ is now a number, nothing stops us from fitting $\yhat = wx + b$
with the closed-form solution from Section~1.1.
Here $\xbar = 3.5$ and $\ybar = 0.5$:

\begin{center}
  \renewcommand{\arraystretch}{1.3}
  \begin{tabular}{ccccc}
    \toprule
    $i$ & $x_i - \xbar$            & $y_i - \ybar$
        & $(x_i-\xbar)(y_i-\ybar)$ & $(x_i-\xbar)^2$                                             \\
    \midrule
    1   & $-2.5$                   & $-0.5$           & $\phantom{-}1.25$  & $6.25$              \\
    2   & $-1.5$                   & $-0.5$           & $\phantom{-}0.75$  & $2.25$              \\
    3   & $-0.5$                   & $\phantom{-}0.5$ & $-0.25$            & $0.25$              \\
    4   & $\phantom{-}0.5$         & $-0.5$           & $-0.25$            & $0.25$              \\
    5   & $\phantom{-}1.5$         & $\phantom{-}0.5$ & $\phantom{-}0.75$  & $2.25$              \\
    6   & $\phantom{-}2.5$         & $\phantom{-}0.5$ & $\phantom{-}1.25$  & $6.25$              \\
    \midrule
        &                          & \textbf{Sum}     & $\sol[1.5cm]{3.5}$ & $\sol[1.5cm]{17.5}$ \\
    \bottomrule
  \end{tabular}
\end{center}

% ── \sol applied: student computes w and b from the sums ─────────────────────
\[
  w = \sol{\dfrac{3.5}{17.5} = 0.2},
  \qquad
  b = \sol{0.5 - 0.2 \times 3.5 = -0.2}.
\]

The fitted line is $\yhat = \sol{0.2\,x - 0.2}$.
% ─────────────────────────────────────────────────────────────────────────────

\textbf{Why not use linear regression?}
As a \emph{classifier}, the line is not terrible: thresholding at
$\yhat = 0.5$ gives the boundary $x = 3.5$, right in the middle of the
overlap.
The trouble starts when we try to read $\yhat$ as a
\emph{probability} of class T (see the accompanying marimo notebook):

\begin{itemize}
  \item \textbf{Values outside $[0, 1]$.}
        At $x = 0$ the line predicts $\yhat = \sol[1.5cm]{-0.2}$, and at $x = 8$ it
        predicts $\yhat = \sol[1.5cm]{1.4}$.
        Neither is a valid probability.
  \item \textbf{Constant slope.}
        Each extra unit of $x$ adds the same $0.2$, whether we are deep inside
        one class or in the uncertain overlap.
        Intuitively, moving from $x = 3$ to $x = 4$ should change our belief a lot,
        while moving from $x = 5$ to $x = 6$ (already confidently T) should change
        it very little.
  \item \textbf{Sensitivity to outliers.}
        A single far-away point, e.g.\ $(x, y) = (20, 1)$, flattens the line
        and drags the $0.5$-threshold to the right, so points that were
        correctly classified before can become misclassified.
\end{itemize}

\subsection{What Shape Do We Want?}

What we really want is a model that outputs a \emph{probability}:
\[
  \hat{p} = P(y = 1 \mid x) \in (0, 1).
\]
For the toy data, a sensible curve should:
\begin{itemize}[nosep]
  \item stay close to $0$ for small $x$ (confidently F),
  \item stay close to $1$ for large $x$ (confidently T),
  \item rise \emph{smoothly} through the overlap region, passing $0.5$ near
        $x = 3.5$, where we are least certain.
\end{itemize}
This is an \textbf{S-shaped} curve (see the marimo notebook), not a
straight line.

% ══════════════════════════════════════════════════════════════════════════════
\section{The Sigmoid Function: Motivation via Log-Odds}

We would still like to use a linear function $wx + b$: it is simple,
interpretable, and we already know how to work with it.
The problem is a mismatch of ranges:
\[
  wx + b \in (-\infty, \infty)
  \qquad \text{but} \qquad
  p \in (0, 1).
\]
So instead of forcing $p$ itself to be linear in $x$, we \textbf{transform}
$p$ into a quantity that lives on all of $\mathbb{R}$, and make
\emph{that} quantity linear.
We do this in two steps.

\textbf{Step 1: probability $\to$ odds.}
The \textbf{odds} of an event with probability $p$ are:
\[
  \text{odds} = \frac{p}{1 - p} \in (0, \infty).
\]
This removes the upper limit of $1$, but the odds are still positive.

\textbf{Step 2: odds $\to$ log-odds.}
Taking the logarithm maps $(0, \infty)$ onto all of $\mathbb{R}$:
\[
  \text{log-odds} = \log \frac{p}{1 - p} \in (-\infty, \infty).
\]
This quantity is also called the \textbf{logit} of $p$.

\begin{center}
  \renewcommand{\arraystretch}{1.3}
  \begin{tabular}{cccc}
    \toprule
    $p$   & odds $= \dfrac{p}{1-p}$ & log-odds             & Interpretation    \\
    \midrule
    $0.1$ & $1/9$                   & $\sol[1.5cm]{-2.20}$ & likely F          \\
    $0.5$ & $\sol[1.5cm]{1}$        & $\sol[1.5cm]{0}$     & completely unsure \\
    $0.9$ & $9$                     & $\sol[1.5cm]{2.20}$  & likely T          \\
    \bottomrule
  \end{tabular}
\end{center}

The log-odds are symmetric around $0$, and $p = 0.5$ maps exactly to $0$.
Now the ranges match, so we model the log-odds as linear in $x$:

\begin{equation}
  \log \frac{p}{1-p} = wx + b.
  \label{eq:logodds}
\end{equation}

\begin{notebox}
  The log-odds is not the only possible choice: any function that maps
  $(0, 1)$ onto $\mathbb{R}$ would work (for example, the inverse of the
  normal CDF gives \emph{probit} regression).
  The log-odds is preferred because it is simple and
  interpretable: increasing $x$ by one unit multiplies the odds by $e^{w}$.
\end{notebox}

\subsection{Solving for \texorpdfstring{$p$}{p}}

Exponentiating both sides of \eqref{eq:logodds} and solving for $p$:

\begin{align*}
  \frac{p}{1-p}     & = e^{wx+b}        \\[4pt]
  p                 & = (1-p)\,e^{wx+b} \\[4pt]
  p\,(1 + e^{wx+b}) & = e^{wx+b}
\end{align*}

% ── \sol applied: student completes the derivation ────────────────────────────
\begin{resultbox}
  \[
    \hat{p} \;=\; \sol[9cm]{\frac{1}{1 + e^{-(wx+b)}} \;=\; \sigmoid(wx+b)}
  \]
\end{resultbox}
% ─────────────────────────────────────────────────────────────────────────────

This function $\sigmoid(z) = \dfrac{1}{1+e^{-z}}$ is called the
\textbf{sigmoid} (or logistic) function.
It is exactly the smooth S-shaped curve we asked for: by making the
log-odds linear, the probability automatically becomes S-shaped in $x$.

\begin{notebox}
  \textbf{Key properties of $\sigmoid(z)$.}
  \begin{itemize}[nosep]
    \item Output is always in $(0, 1)$ — valid as a probability.
    \item $\sigmoid(0) = 0.5$: the decision boundary is where $wx + b = 0$.
    \item As $z \to +\infty$, $\sigmoid(z) \to 1$; as $z \to -\infty$,
          $\sigmoid(z) \to 0$.
    \item Its derivative has the elegant form
          $\sigmoid'(z) = \sigmoid(z)\bigl(1 - \sigmoid(z)\bigr)$
          (see Appendix~\ref{app:sigmoid-deriv}).
  \end{itemize}
\end{notebox}

% ══════════════════════════════════════════════════════════════════════════════
\section{The Logistic Regression Model}

So far we have used a single feature $x$, and the sigmoid's input is the
number $z = wx + b$.
With more features, \textbf{only $z$ changes}: it is still a single number,
it just collects more terms.

\subsection{Two Features: Writing It Out}

Suppose each observation has two features $x_1$ and $x_2$
(e.g.\ hours studied and hours slept).
The log-odds are now linear in both:
\[
  z = b + w_1 x_1 + w_2 x_2,
  \qquad
  \hat{p} = \sigmoid(z) = \frac{1}{1 + e^{-z}}.
\]

To write this compactly, collect the parameters and the features
(with a leading $1$ for the bias, as in Section~1.1) into column vectors:
\[
  \bth = \begin{bmatrix} b \\ w_1 \\ w_2 \end{bmatrix},
  \qquad
  \bx = \begin{bmatrix} 1 \\ x_1 \\ x_2 \end{bmatrix}.
\]
Then $\bth^{\top}$ is a \emph{row} vector, and multiplying a row by a column
gives back exactly the same sum:
\[
  \bth^{\top}\bx
  = \begin{bmatrix} b & w_1 & w_2 \end{bmatrix}
  \begin{bmatrix} 1 \\ x_1 \\ x_2 \end{bmatrix}
  = \sol[5cm]{b \cdot 1 + w_1 x_1 + w_2 x_2}.
\]
Size check: $(1 \times 3)\,(3 \times 1) = (1 \times 1)$, a \textbf{single number}.

\begin{example}[Computing $\hat{p}$ with two features]
  Let $\bth = [-3,\; 1,\; 0.5]^{\top}$ and consider a student with
  $x_1 = 2$ hours studied and $x_2 = 4$ hours slept, so
  $\bx = [1,\; 2,\; 4]^{\top}$.

  % ── \sol applied: student evaluates z and p ────────────────────────────────
  \[
    z = \bth^{\top}\bx
    = \begin{bmatrix} -3 & 1 & 0.5 \end{bmatrix}
    \begin{bmatrix} 1 \\ 2 \\ 4 \end{bmatrix}
    = \sol[6cm]{(-3)(1) + (1)(2) + (0.5)(4) = 1}
  \]
  \[
    \hat{p} = \sigmoid(z) = \sol[6cm]{\frac{1}{1 + e^{-1}} \approx 0.73}
  \]
  % ──────────────────────────────────────────────────────────────────────────

  Since $\hat{p} \ge 0.5$, the predicted class is $\hat{y} = \sol[1cm]{1}$ (pass).
\end{example}

\subsection{Any Number of Features}

For $d$ features the pattern is the same; the vectors just get longer:
\[
  z = \bth^{\top}\bx
  = \underbrace{b + w_1 x_1 + w_2 x_2 + \cdots + w_d x_d}_{\text{still one number}}.
\]

\begin{notebox}
  No matter how many features we have, $\bth^{\top}\bx$ is a
  \textbf{single number} (a scalar).
  The sigmoid only ever sees one number $z$, so everything we learned
  about $\sigmoid(z)$ for one feature carries over unchanged.
  The decision boundary $z = 0$ is a point when $d = 1$, a line when $d = 2$,
  and a plane (hyperplane) for larger $d$.
\end{notebox}

Writing the observation index $i$ explicitly, the model is:

\begin{equation}
  \hat{p}^{(i)} = \sigmoid\!\left(\bth^{\top}\bx^{(i)}\right)
  = \frac{1}{1 + e^{-\bth^{\top}\bx^{(i)}}},
  \label{eq:logreg}
\end{equation}

where $\bth = [b,\, w_1,\, \ldots,\, w_d]^{\top}$ and $\bx^{(i)}$ includes a
leading 1 for the bias (same design-matrix convention as Section~1.1).

The predicted class is:
\[
  \hat{y}^{(i)} =
  \begin{cases} 1 & \text{if } \hat{p}^{(i)} \ge 0.5 \\ 0 & \text{otherwise.} \end{cases}
\]

% ══════════════════════════════════════════════════════════════════════════════
\section{Cost Function: Cross-Entropy via MLE}

\begin{notebox}
  \textbf{Why not MSE again?}
  Technically, we \emph{could} reuse the MSE from Section~1.1,
  $J(\bth) = \frac{1}{n}\sum_{i}\bigl(y^{(i)} - \hat{p}^{(i)}\bigr)^2$,
  but it is a poor choice for classification:
  \begin{enumerate}[nosep]
    \item \textbf{Weak learning from big mistakes.}
          When the model is confidently wrong ($\hat{p}$ near $0$ or $1$),
          the sigmoid is flat, so the MSE gradient is tiny and the model
          barely corrects itself.
    \item \textbf{Non-convex.}
          Combined with the sigmoid, MSE is non-convex in $\bth$, so gradient
          descent can stall on plateaus or local minima.
    \item \textbf{Wrong noise model.}
          MSE suits continuous outputs with Gaussian noise; our labels are
          $0$/$1$, which are naturally modelled as Bernoulli.
  \end{enumerate}
  Deriving the loss from the Bernoulli likelihood (below) fixes all three.
\end{notebox}

\subsection{Setting Up the Likelihood}

We model each label $y^{(i)}$ as a Bernoulli random variable with success
probability $\hat{p}^{(i)}$:
\[
  P\!\left(y^{(i)} \mid \bx^{(i)};\,\bth\right)
  = \left(\hat{p}^{(i)}\right)^{y^{(i)}}
  \left(1 - \hat{p}^{(i)}\right)^{1 - y^{(i)}}.
\]

\textbf{Reading the formula.}
The model gives $\hat{p}^{(i)}$, the probability that $y^{(i)} = 1$.
What we want is the probability of the label we \emph{actually observed}:
\[
  P\!\left(y^{(i)} \mid \bx^{(i)};\,\bth\right) =
  \begin{cases}
    \hat{p}^{(i)}     & \text{if } y^{(i)} = 1, \\[2pt]
    1 - \hat{p}^{(i)} & \text{if } y^{(i)} = 0.
  \end{cases}
\]
The Bernoulli formula is just a compact way to write these two cases in one
line. The exponents work like a switch:

% ── \sol applied: student checks both cases of the switch ────────────────────
\begin{center}
  \renewcommand{\arraystretch}{1.4}
  \begin{tabular}{lll}
    $y = 1$: & $\hat{p}^{\,1}\,(1-\hat{p})^{0}$
             & $= \sol[2.5cm]{\hat{p} \cdot 1 = \hat{p}}$         \\
    $y = 0$: & $\hat{p}^{\,0}\,(1-\hat{p})^{1}$
             & $= \sol[2.5cm]{1 \cdot (1-\hat{p}) = 1 - \hat{p}}$
  \end{tabular}
\end{center}
% ─────────────────────────────────────────────────────────────────────────────

Anything raised to the power $0$ equals $1$, so the factor that does not
match the observed label simply disappears.

\begin{example}[Probability of the observed labels]
  A model gives the following predictions for four observations:

  % ── \sol applied: student fills in the probability of each observed label ──
  \begin{center}
    \renewcommand{\arraystretch}{1.4}
    \begin{tabular}{cccc}
      \toprule
      $i$ & Observed $y^{(i)}$                                                             & Predicted $\hat{p}^{(i)}$
          & $\left(\hat{p}^{(i)}\right)^{y^{(i)}}\left(1-\hat{p}^{(i)}\right)^{1-y^{(i)}}$                                                                         \\
      \midrule
      1   & 0                                                                              & 0.1                       & $0.1^{0} \times 0.9^{1} = 0.9$            \\
      2   & 1                                                                              & 0.8                       & $0.8^{1} \times 0.2^{0} = 0.8$            \\
      3   & 1                                                                              & 0.6                       & $\sol[4cm]{0.6^{1} \times 0.4^{0} = 0.6}$ \\
      4   & 0                                                                              & 0.3                       & $\sol[4cm]{0.3^{0} \times 0.7^{1} = 0.7}$ \\
      \bottomrule
    \end{tabular}
  \end{center}

  Each entry is the probability the model assigns to what \emph{really}
  happened. A value close to $1$ means the model predicted that observation
  well; a value close to $0$ means it was confidently wrong.

  Assuming the observations are independent, the probability of seeing
  \emph{all four} labels is the product:
  \[
    \mathcal{L}
    = \sol[8cm]{0.9 \times 0.8 \times 0.6 \times 0.7 \approx 0.302}
  \]

  Now compare a worse model that predicts $\hat{p} = 0.6,\ 0.3,\ 0.4,\ 0.8$
  for the same four observations:
  \[
    \mathcal{L}
    = \sol[8cm]{0.4 \times 0.3 \times 0.4 \times 0.2 \approx 0.010}
  \]
  The better model gives the observed data a much higher probability.
  \textbf{Maximum likelihood estimation (MLE)} simply chooses the
  parameters $\bth$ that make this product as large as possible.
\end{example}

In general, assuming the $n$ observations are independent, the joint
likelihood is:
\[
  \mathcal{L}(\bth)
  = \prod_{i=1}^{n}
  \left(\hat{p}^{(i)}\right)^{y^{(i)}}
  \left(1 - \hat{p}^{(i)}\right)^{1 - y^{(i)}}.
\]

\subsection{From MLE to the Cross-Entropy Loss}

Maximising $\mathcal{L}$ is equivalent to maximising the log-likelihood
$\ell = \log \mathcal{L}$ (log is monotone):

% ── \sol applied: student expands the log-likelihood ─────────────────────────
\[
  \ell(\bth)
  = \sol[10cm]{\sum_{i=1}^{n}
    \Bigl[y^{(i)} \log \hat{p}^{(i)}
      + (1 - y^{(i)}) \log\!\left(1 - \hat{p}^{(i)}\right)\Bigr].}
\]
% ─────────────────────────────────────────────────────────────────────────────

Minimising the \textbf{negative} log-likelihood (and averaging over $n$) gives
the \textbf{Binary Cross-Entropy (BCE)} loss:

\begin{resultbox}
  \[
    J(\bth)
    = \sol[10cm]{-\frac{1}{n}\sum_{i=1}^{n}
      \Bigl[y^{(i)} \log \hat{p}^{(i)}
        + (1 - y^{(i)}) \log\!\left(1 - \hat{p}^{(i)}\right)\Bigr]}
  \]
\end{resultbox}

\begin{notebox}
  \textbf{Why cross-entropy?}
  When $y^{(i)} = 1$, only the first term $-\log\hat{p}^{(i)}$ matters:
  we are penalised heavily if $\hat{p}^{(i)}$ is close to 0 (confident and
  wrong).  When $y^{(i)} = 0$, only the second term $-\log(1-\hat{p}^{(i)})$
  matters.  The two terms switch on and off depending on the true label.
\end{notebox}

% ══════════════════════════════════════════════════════════════════════════════
\section{The Gradient and Why We Need Gradient Descent}

Unlike MSE in linear regression, there is \textbf{no closed-form solution}
for the minimiser of $J(\bth)$.  The cross-entropy loss is convex (it has a
unique global minimum), but we cannot invert the normal equations here because
the sigmoid is non-linear.

Taking the gradient with respect to $\bth$ (using the chain rule and the
identity $\sigmoid'(z) = \sigmoid(z)(1-\sigmoid(z))$):

\begin{resultbox}
  \[
    \nabla_{\bth}\,J
    = \sol[8cm]{\frac{1}{n}\sum_{i=1}^{n}
      \Bigl(\hat{p}^{(i)} - y^{(i)}\Bigr)\,\bx^{(i)}}
  \]
\end{resultbox}

\begin{notebox}
  \textbf{Remarkably clean form.}
  The gradient has exactly the same structure as in linear regression —
  it is a weighted sum of the input vectors, where the weights are the
  \emph{prediction errors} $(\hat{p}^{(i)} - y^{(i)})$.
  The sigmoid ``absorbed'' all the complexity through the chain rule.
\end{notebox}

Because there is no closed form, we use \textbf{gradient descent}: starting
from an initial $\bth^{(0)}$, we repeatedly update
\[
  \bth \;\leftarrow\; \bth - \alpha\,\nabla_{\bth}\,J,
\]
where $\alpha > 0$ is the \textbf{learning rate}.
Gradient descent will be covered in detail in the next section.

% ══════════════════════════════════════════════════════════════════════════════
\section{Summary}

\begin{center}
  \renewcommand{\arraystretch}{1.5}
  \begin{tabular}{ll}
    \toprule
    \textbf{Model}           & $\hat{p} = \sigmoid(\bth^{\top}\bx)$                                        \\
    \textbf{Loss function}   & BCE: $J = -\dfrac{1}{n}\sum\bigl[y\log\hat{p} + (1-y)\log(1-\hat{p})\bigr]$ \\
    \textbf{Gradient}        & $\nabla_{\bth}J = \dfrac{1}{n}\sum(\hat{p}^{(i)} - y^{(i)})\,\bx^{(i)}$     \\
    \textbf{Solution method} & Gradient descent (no closed form)                                           \\
    \textbf{Works well when} & The decision boundary is (approximately) linear                             \\
    \bottomrule
  \end{tabular}
\end{center}

\bigskip
\textbf{Up next (Section~1.3):} What if no linear boundary can separate the
classes?  We revisit the XOR dataset to show that both linear regression and
logistic regression fail — motivating the need for non-linearity and neural
networks.

% ══════════════════════════════════════════════════════════════════════════════
\appendix
\section{Derivative of the Sigmoid Function}
\label{app:sigmoid-deriv}

Let $\sigmoid(z) = (1 + e^{-z})^{-1}$.  By the chain rule:

\begin{align*}
  \sigmoid'(z)
   & = -\frac{-e^{-z}}{(1+e^{-z})^2}
  = \frac{e^{-z}}{(1+e^{-z})^2}.
\end{align*}

Factor out $\dfrac{1}{1+e^{-z}}$ from the numerator and denominator:

\[
  \sigmoid'(z)
  = \frac{1}{1+e^{-z}} \cdot \frac{e^{-z}}{1+e^{-z}}
  = \sigmoid(z) \cdot \frac{1+e^{-z}-1}{1+e^{-z}}
  = \sigmoid(z)\bigl(1 - \sigmoid(z)\bigr).
\]

\begin{resultbox}
  \[
    \sigmoid'(z) = \sigmoid(z)\bigl(1 - \sigmoid(z)\bigr)
  \]
\end{resultbox}

This compact form makes the chain-rule calculations in gradient descent
particularly clean.

\end{document}
